% Replacement text for Sections 1--3 of claude-version_merged_EN.tex.
% All revised material is red for review.
\begingroup\color{red}
\hypersetup{citecolor=red,linkcolor=red}

%======================================================================
\section{Introduction}
\label{sec:introduction}

Epidemics spread at more than one scale. Inside a community, infection may pass through a detailed network of contacts between people \citep{keeling2005,pastorsatorras2015,kiss2017}. Between communities, the available data are often much less detailed. For example, they may give only passenger numbers, commuting flows, or a weighted link between two places \citep{balcan2009,barbosa2018human}. We therefore need a model that can use a detailed network inside each community and simpler data between communities.

This situation is common in practice. A local contact network may be observed, reconstructed, or generated to match known properties such as clustering and mean shortest-path length \citep{eubank2004,danon2011}. In contrast, links between communities may be known only through transport, migration, commuting, or mobile-phone data \citep{barbosa2018human,chang2021}. Such data can tell us the total strength of connection between two communities, but not which people form the connection. We may therefore know the total coupling without knowing the bridge nodes or their infection states.

If every contact is known, a continuous-time Markov chain (CTMC) can describe the epidemic, and a Gillespie-type method can simulate it exactly \citep{gillespie1977exact,gibson2000}. Efficient methods are also available for temporal, large, and sparse networks \citep{vestergaard2015,kuryliak2025simulating}. However, these methods still need the individual contacts that may transmit infection. Metapopulation models use aggregate flows between communities, but usually also replace the detailed local network by population counts \citep{colizza2007reaction,colizza2008,balcan2010modeling,ajelli2010}. Our problem lies between these two cases: we want to keep node-level stochastic spread inside every community while using only aggregate information between communities.

The basic two-scale simulator was presented in an earlier workshop paper \citep{kuryliak2025towards}. It runs one stochastic local simulator for each community. A separate macro layer uses the current community states and the aggregate coupling to generate infections that enter from another community. The present article gives a full description of this method and studies its approximation error. In particular, we ask what is lost when real cross-community edges and their endpoints are replaced by an aggregate connection, and what is lost when an imported infection is placed on a selected susceptible node.

Our main question is: when does this aggregate interface reproduce the behaviour of the full-network CTMC well enough for the quantity of interest? We study arrival times, their distributions, and the physical validity of the transmission route. In particular, similar arrival times do not prove that the route was physically possible: the aggregate model can export infection before any real bridge node has been infected. The complete microscopic network is known in our controlled experiments, so it can be used as a reference. The two-scale simulator receives only the local networks and aggregate connections. These experiments test the method; they are not a calibration of one particular epidemic.

This article makes four contributions.
\begin{enumerate}
\item We give a complete description of the two-scale simulator, including the time-grid scheduler used in the experiments (Section~\ref{sec:method}, Algorithm~\ref{alg:scheduler}).
\item We show exactly what the ideal continuous-time method simulates. It is exact for a \emph{surrogate network}: the real cross-community edges are replaced by weak all-to-all connections with the same total weight (Proposition~\ref{prop:surrogate}). It is not, in general, exact for the original physical network. We then identify the resulting sources of error and analyse a single export (Proposition~\ref{prop:ordering}).
\item We compare the method with full-network Gillespie simulation for single exports and chains of communities. Permutation baselines show how much difference can arise from finite sampling alone (Sections~\ref{sec:design} and~\ref{sec:results}).
\item We give an empirical indicator of accuracy for our benchmarks. It compares the spread of infection times inside a community with the typical time needed to cross one interface edge (Section~\ref{sec:res_descriptor}).
\end{enumerate}

We consider susceptible--infected (SI) dynamics on static synthetic networks and use one fixed synchronisation step. Section~\ref{sec:related} explains the relation to earlier methods. Section~\ref{sec:method} describes the simulator. Section~\ref{sec:analysis} explains what the aggregate interface preserves and loses. Sections~\ref{sec:design} and~\ref{sec:results} present the experiments and results. Section~\ref{sec:discussion} discusses use, cost, and limitations.

\clearpage

%======================================================================
\section{Related work}
\label{sec:related}

\paragraph{Simulation of stochastic events.}
Gillespie's direct method draws the next event and its time from the current event rates \citep{gillespie1977exact}. The next reaction method keeps a separate clock for each possible event \citep{gibson2000}. Both methods are exact for the CTMC they simulate. When an event rate changes with time, its rate can be accumulated until it reaches a random threshold \citep{ethier1986}. \citet{anderson2007} used this idea for systems with time-dependent rates and delays. Our macro clock uses the same idea, but its rate changes because the local epidemic states change. Network simulators can save work by updating only the rates affected by the latest event \citep{vestergaard2015,kuryliak2025simulating}. Our local simulators use the implementation of \citet{kuryliak2025simulating}.

Some multiscale methods speed up chemical-reaction models by approximating reactions that occur much faster than others \citep{haseltine2002,cao2005slow}. Our method removes a different kind of detail. It keeps the local infection events, but does not know which individual cross-community edges are currently able to transmit. Even so, the relative speeds of local spread and cross-community transmission help explain its error.

\paragraph{Models with several communities.}
Metapopulation models represent each subpopulation by counts and connect the subpopulations by mobility flows. They have been used to study global spread and the conditions for invasion \citep{colizza2007reaction,colizza2008,balcan2009,balcan2010modeling}. Detailed mobility data can also be combined with local transmission models \citep{chang2021}. Agent-based models instead represent individuals throughout the model \citep{eubank2004,kerr2021}. \citet{ajelli2010} compared these two broad approaches. The method studied here keeps individual nodes and contacts inside each community, but replaces the unknown individual contacts between communities by one aggregate stochastic interface.

\paragraph{Two levels of mixing and co-simulation.}
Some epidemic models combine contacts in small local groups with uniform mixing at a wider level \citep{ball1997}. Our surrogate network is related to this idea, but uniform mixing is used only between connected pairs of communities, and its strength is set by their aggregate weight.

The implementation also has features of co-simulation. In co-simulation, separate subsystem solvers move forward on a shared time grid and exchange a limited amount of information \citep{gomes2018cosim}. Here, the local community simulators are those subsystems. They send their infected and susceptible counts to the macro layer. In this implementation, those exchanged counts are used to draw a cross-community infection time and then to choose its receiving community and target node. The scheduler can therefore treat the local simulators as black boxes.

\clearpage

%======================================================================
\section{The two-scale simulation method}
\label{sec:method}

We first describe the full network used as the reference. We then describe the information given to the two-scale simulator and how it generates an imported infection. Two versions must be kept separate. The \emph{ideal} version works in continuous time and is exact for the surrogate network defined later, not generally for the original physical network. The implemented version checks the local simulators on a time grid and therefore has an additional numerical approximation.

\subsection{Full-network reference model}
\label{sec:reference}

Let $G=(V,E)$ be the contact network. Its nodes are divided into $M$ non-overlapping communities:
\begin{equation}
\label{eq:partition}
V=V_1\cup V_2\cup\cdots\cup V_M,
\qquad N_i=|V_i|.
\end{equation}
Here, the sets $V_i$ do not overlap, and $N_i$ is the number of nodes in community $i$. We write $E_{ii}$ for contacts inside community $i$, and
\[
E_{ij}=\{(u,v)\in E:u\in V_i,\,v\in V_j\},
\qquad i\neq j,
\]
for contacts from community $i$ to community $j$. We treat each contact as an oriented transmission channel. An undirected contact therefore gives two channels, one in each direction.

Every node is susceptible ($S$), infected ($I$), or removed ($R$). Its full state at time $t$ is $X(t)$, and $X_i(t)$ is the part of this state in community $V_i$. In community $i$, let $S_i(t)$ and $I_i(t)$ be the sets of susceptible and infected nodes. Their sizes are $s_i(t)$ and $y_i(t)$, and their fractions are $\sbar_i(t)=s_i(t)/N_i$ and $\ibar_i(t)=y_i(t)/N_i$. For a state $x$, the notation $x^{v:I}$ means the same state after node $v$ becomes infected, and $x^{u:R}$ means the state after node $u$ is removed.

The full model follows two simple rules. A susceptible node $v$ receives infection from each infected neighbour $u$ at rate $\beta_{uv}$. An infected node is removed at rate $\gamma$. We can write both rules as
\begin{equation}
\label{eq:full_generator}
q_v(x)=\sum_{\substack{u:x_u=I\\(u,v)\in E}}\beta_{uv}
\quad\text{for }x_v=S,
\qquad
r_u(x)=\gamma
\quad\text{for }x_u=I.
\end{equation}
Thus, $q_v(x)$ is the total infection rate of a susceptible node $v$, and $r_u(x)$ is the removal rate of an infected node $u$; the corresponding rate is zero for nodes in other states. A rate describes how often an event occurs: for a constant rate $q$, the mean waiting time is $1/q$. These rates also define the CTMC generator $\mathcal{L}_G$ used later; no additional epidemic rule is hidden in that notation. For a cross-community channel, we write $\beta_{uv}=\beta_{\mathrm{out}}w_{uv}$, where $w_{uv}$ is its contact weight. We call exact simulation of these full-network rates \Full{}. All experiments in this article use SI dynamics, so $\gamma=0$.

\subsection{Aggregate interface}
\label{sec:interface}

The two-scale simulator, denoted \TS{}, knows every local graph $G_i=(V_i,E_{ii})$ and its rates. It does not know the individual cross-community channels in $E_{ij}$. For each ordered pair of communities, it knows only the total weight
\begin{equation}
\label{eq:weight}
W_{ij}=\sum_{(u,v)\in E_{ij}}w_{uv},
\qquad i\neq j.
\end{equation}
If all weights equal one, $W_{ij}$ is the number of directed channels from $V_i$ to $V_j$. In an application, this total may instead come from transport, migration, mobility, or contact data. The macro layer also receives the current susceptible and infected counts from each local simulator. It does not know which nodes are endpoints of real cross-community edges. Table~\ref{tab:information} gives the main difference, and Figure~\ref{fig:schematic} shows it visually.

\begin{table}[tbp]
\color{red}
\centering
\caption{\textcolor{red}{Information used by the two models. The ``target rule'' in \TS{} is uniform among susceptible nodes in all experiments in this article.}}
\label{tab:information}
\small
\begin{tabular}{>{\raggedright\arraybackslash}p{0.42\linewidth}>{\raggedright\arraybackslash}p{0.23\linewidth}>{\raggedright\arraybackslash}p{0.25\linewidth}}
\toprule
Information or action & \Full{} & \TS{}\\
\midrule
Local contacts, rates, and node states & known & known by local simulators\\
Individual cross-community edges and bridge nodes & known & unknown\\
Total cross-community weight $W_{ij}$ & obtained from the edges & given directly\\
Cross-community event rate & uses states of real edge endpoints & uses $\ibar_i$, $\sbar_j$, and $W_{ij}$\\
Target of a cross-community infection & endpoint of the active edge & a susceptible node drawn by the target rule\\
Number of possible SI infections through one bundle & at most the number of its channels & at most $s_j$ target infections\\
\bottomrule
\end{tabular}
\end{table}

\begin{figure}[tbp]
\color{red}
\centering
\includegraphics[width=\linewidth]{Images/twoscale/fig_method_schematic.pdf}
\caption{\textcolor{red}{Three views of the interface. (a) In the full network, infection can cross only through the shown physical edges and their bridge nodes. (b) The two-scale simulator keeps a separate local simulator for each community. The macro layer receives community counts, uses $W_{ij}$ to generate a cross-community infection, and places it on a uniformly selected susceptible node. (c) The ideal two-scale process is equivalent to a surrogate network in which every ordered node pair across the two communities has a weak connection. The connections together have total weight $W_{ij}$. Orange nodes are infected and blue nodes are susceptible.}}
\label{fig:schematic}
\end{figure}

\subsection{Local simulators}
\label{sec:local}

Each community has a separate Gillespie simulator for its local graph $G_i$. It applies the infection and removal rates in~\eqref{eq:full_generator} using only the local edges $E_{ii}$. The implementation updates only rates affected by the latest event, following \citet{kuryliak2025simulating}. Each local simulator has its own random-number stream.

The macro layer can ask a local simulator to move to a specified time. It can also import an infection by changing one selected susceptible node to $I$. An import changes the local rates. The corrected implementation keeps the unused part of the exponential threshold for the next local event and recalculates its waiting time using the new total rate. It then selects the event type and node using the updated rates. A new local threshold is drawn only if there is no existing one. This follows the same accumulated-rate principle explained in Section~\ref{sec:timing}. For prescribed import times and target nodes, the local simulator follows the CTMC rules of its own network. The approximation comes from the way the macro layer generates and places imports.

\subsection{Macro layer}
\label{sec:macro}

The macro layer represents each community by its current fractions and represents each directed connection by $W_{ij}$. Its rate for an infection from community $i$ to community $j$ is
\begin{equation}
\label{eq:macro_rate}
\lambda_{ij}(t)=\beta_{\mathrm{out}}W_{ij}\ibar_i(t)\sbar_j(t),
\qquad
\Theta(t)=\sum_{i\neq j}\lambda_{ij}(t).
\end{equation}
The factors have a direct meaning: $\beta_{\mathrm{out}}W_{ij}$ gives the total connection strength, $\ibar_i$ is the infected fraction in the source, and $\sbar_j$ is the susceptible fraction in the receiving community. The sum $\Theta$ is the rate of any cross-community infection.

When such an event occurs, the pair $(i,j)$ is selected with probability $\lambda_{ij}/\Theta$. Write $\Pi_{ij}(v\mid x)$ for the probability of selecting target node $v$ when the current state is $x$. The default rule chooses uniformly from the susceptible set $S_j$, so $\Pi_{ij}(v\mid x)=1/s_j$ for every susceptible node $v\in S_j$. Other target rules can be used if bridge-node or boundary information is available. The optional calibration factor in the implementation is one in all experiments.

This aggregate rate is generally an approximation to the rate across the real physical edges. For fixed community counts, it is the expected rate if the endpoints of those edges are placed uniformly at random, before their identities, states, and epidemic history are known. It becomes the actual rate for the uniform all-to-all surrogate interface described in Section~\ref{sec:analysis}. The rate $\Theta$ changes after any local event or imported infection that changes a community count.

\subsection{Timing and synchronisation}
\label{sec:timing}

After a macro event at time $t_0$, draw a random threshold $E\sim\Exp(1)$: an exponential random number with mean one. Starting from zero, accumulate the total macro rate $\Theta$. The next macro event occurs at the first time $t^\ast$ when
\begin{equation}
\label{eq:hazard_integral}
\int_{t_0}^{t^\ast}\Theta(s)\,ds=E.
\end{equation}
In words, the area under the rate curve grows until it first reaches the random threshold. If it never reaches the threshold, no further macro event occurs. This is the standard cumulative-rate construction \citep{ethier1986,anderson2007}. If $\Theta$ is constant, it gives the familiar Gillespie waiting time $(t^\ast-t_0)=-\ln U/\Theta$, where $U$ is uniform on $(0,1)$. Exact evaluation of this first crossing, together with exact local simulators, defines the ideal two-scale process. Section~\ref{sec:surrogate} shows that this ideal process is exact for the surrogate network.

The implemented scheduler checks progress on intervals of length $\tau$. Let $A$ be the amount of macro rate already accumulated since the previous macro event. For the current interval $[t,t_g)$, it copies every local simulator, including its random-number state, and advances each copy to the midpoint. The copies give a midpoint estimate $\Theta_{\mathrm{mid}}$. If the estimated accumulation remains below $E$, the real simulators move to $t_g$. If it reaches $E$, the scheduler estimates the crossing time by
\begin{equation}
\label{eq:interpolation}
t^\ast=t+\frac{E-A}{\Theta_{\mathrm{mid}}}.
\end{equation}
The numerator is the amount left before the threshold, and the denominator is the estimated rate. The scheduler then advances the real simulators to $t^\ast$, draws the source and receiving communities, selects the target node, and imports the infection. Algorithm~\ref{alg:scheduler} gives the full sequence.

\begin{algorithm}[H]
\color{red}
\caption{\textcolor{red}{Grid-synchronised two-scale simulation of one realisation.}}
\label{alg:scheduler}
\small
\fbox{\begin{minipage}{0.94\linewidth}
\textbf{Input:} local graphs $G_i$, weights $W_{ij}$, rates, step $\tau$, horizon $T_{\mathrm{end}}$, and initial state.
\begin{enumerate}
\item Set $t=0$ and accumulated rate $A=0$; draw $E\sim\Exp(1)$.
\item While $t<T_{\mathrm{end}}$ and infection can still spread, set $t_g=\min(t+\tau,T_{\mathrm{end}})$ and process the interval up to $t_g$.
\item Copy every local simulator, including its random-number state. Advance the copies to the midpoint of the remaining interval and calculate $\Theta_{\mathrm{mid}}$ from~\eqref{eq:macro_rate}.
\item If $A+\Theta_{\mathrm{mid}}(t_g-t)<E$, advance the real simulators to $t_g$, add this amount to $A$, set $t=t_g$, and return to step 2.
\item Otherwise, calculate $t^\ast$ from~\eqref{eq:interpolation} and advance the real simulators to $t^\ast$.
\item Select $(i,j)$ with probability $\lambda_{ij}(t^\ast)/\Theta(t^\ast)$. Select $v$ uniformly from $S_j(t^\ast)$ and import infection at $v$. Update the receiving local simulator, keeping its unused exponential threshold and selecting the local event from the new rates.
\item Set $t=t^\ast$ and $A=0$, draw a new $E\sim\Exp(1)$, and process the rest of the current interval in the same way.
\end{enumerate}
\end{minipage}}
\end{algorithm}

The copied random-number states are important. Unless an import occurs first, a copied simulator follows the same path that the real simulator will follow to the midpoint. The midpoint estimate is exact on an interval with no local event. Its timing error and its look-ahead of at most $\tau/2$ decrease to zero as $\tau$ decreases.

The time grid is needed because the scheduler treats local simulators as black boxes. If they also report their next local-event times, the scheduler can merge those times and accumulate the piecewise-constant value of $\Theta$ exactly. That event-driven version removes the numerical synchronisation approximation. The experiments in this article use the implemented grid version with $\tau=1$.

\subsection{Computational structure and intended use}
\label{sec:cost}

\Full{} stores all local and cross-community channels, and its work per event grows with node degree \citep{kuryliak2025simulating}. \TS{} stores the local graphs and an $M\times M$ matrix of aggregate weights. This is the memory saving described in the workshop paper \citep{kuryliak2025towards}. It can matter when a mobility flow stands for many unknown individual contacts. Local simulators can also use different community-specific settings.

The grid scheduler has its own cost because it copies local states at every time step. Therefore, the method is intended for cases where individual cross-community contacts are unavailable or too numerous to store. It is not expected to be faster for every known, small network. Section~\ref{sec:disc_cost} reports the measured run times; the main evaluation in this article concerns approximation error.

\clearpage
\endgroup
